\section{Sharp density geometry and quantitative polarization}
\label{sec:r12-polarization-foundation}
This appendix reproduces the analytic inputs used by the new evolution
arguments. They are proved here, including the infinite-dimensional
compactness step. No radiality, positivity of the kinetic semigroup,
or real-sector nondegeneracy is used.

Let $\mathcal H$ be any nonzero separable complex Hilbert space.
Finite-dimensional $\mathbb C^d$ and $\ell^2$ are both allowed.
Fix $N\ge1$, $m>1$, $s=m/2$, and a real-even, positive elliptic,
$m$-homogeneous symbol $a$, of class $C^{N+2}$ away from zero.
For $U\in H^s(\mathbb R^N;\mathcal H)$ put
\[
 \rho_U=\|U\|_{\mathcal H}^2,\quad
 M(U)=\int\rho_U,\quad
 A(U)=\int a(\xi)\|\widehat U(\xi)\|_{\mathcal H}^2\,d\xi.
\]
We use either of the following interactions:
\begin{align*}
 &\kappa=1+m/N,
 &&\mathcal F(U)=\int\rho_U^\kappa,
 &&V_\rho=\rho^{\kappa-1};\\
 &\kappa=2,
 &&\mathcal F(U)=\mathcal D(\rho_U,\rho_U),
 &&V_\rho=K*\rho.
\end{align*}
In the second case $1<m<N$, the real even kernel $K$ is homogeneous
of degree $-m$, its interaction form $\mathcal D$ is positive definite,
and
\[
 |K(x)|+|x||\nabla K(x)|\le C|x|^{-m}.
\]
As before, positivity of the Fourier transform of $K$ almost everywhere
suffices for positive definiteness. Pointwise positivity is not needed.
The scalar sharp inequality defines $M_c$ by
\begin{equation}\label{eq:r11-vectorGN}
 \mathcal F(u)\le\kappa
       \left(\frac{M(u)}{M_c}\right)^{\kappa-1}A(u).
\end{equation}

\begin{theorem}[Sharp interpolation and dimension-independent compactness]
\label{thm:r11-vectorcompact}
Inequality \eqref{eq:r11-vectorGN} has the same sharp constant for
$\mathcal H$-valued functions. Every optimizer has the form
\[
 U(x)=\zeta Q(x),\qquad \|\zeta\|_{\mathcal H}=1,
\]
where $Q$ is a real scalar optimizer, up to a constant phase.
If $M(U_n)\to M_c$, $A(U_n)\to1$, and
$\mathcal F(U_n)\to\kappa$, then the sequence is strongly compact
in $H^s$ after translations and internal unitaries. This remains
true when the internal dimensions vary: embed all spaces in
$\ell^2$ before applying the unitaries.
\end{theorem}

\subsection{Sharp interpolation and complete polarization}
The local interaction is the $\kappa$-th power of the $L^\kappa$
density norm. The Hartree interaction is the square of its positive
definite Hilbert density norm. Thus both have the form
$\mathcal F(U)=\|\rho_U\|_{\mathcal B}^\kappa$ for a strictly convex
real density space $\mathcal B$.
In an internal orthonormal basis $U=(u_j)$, the triangle inequality,
scalar sharp interpolation, and sequence H\"older give
\begin{align*}
 \mathcal F(U)^{1/\kappa}
 &\le\sum_j\mathcal F(u_j)^{1/\kappa}\\
 &\le\kappa^{1/\kappa}M_c^{-(\kappa-1)/\kappa}
   \sum_j A(u_j)^{1/\kappa}M(u_j)^{(\kappa-1)/\kappa}\\
 &\le\kappa^{1/\kappa}M_c^{-(\kappa-1)/\kappa}
              A(U)^{1/\kappa}M(U)^{(\kappa-1)/\kappa}.
\end{align*}
Finite partial sums converge in the density norm because the same
bound applied to their tails tends to zero. A scalar optimizer placed
in one component proves sharpness.

Split the components further into real scalar functions $f_j$.
Equality forces every active $f_j$ to be a scalar optimizer, all
$f_j^2$ to be proportional, and all kinetic-to-mass ratios to agree.
The same conclusions hold after any real orthogonal rotation of a
pair $f_i,f_j$. If $f_j/f_i$ took both signs of a nonzero constant,
the squared sum and difference could not be proportional; when the
constant has magnitude one their nonzero supports would be disjoint.
Thus each $f_j$ is one real constant multiple of a fixed scalar
function. This proves complete polarization, including for countably
many components.

\subsection{Compactness modulo internal unitaries}
Ordinary local compactness does not hold directly in an infinite
internal dimension: a fixed spatial shape can move through mutually
orthogonal internal vectors. The following argument removes precisely
that loss of compactness.

\begin{lemma}[Internal-unitary local compactness]
\label{lem:r11-internalcompact}
Let $X=H^s(\mathbb R^N;\mathbb C)$ with $s>0$. From every bounded
sequence in $X\widehat\otimes\ell^2$, one can extract a subsequence
and internal unitaries $W_n$ such that $W_nU_n$ converges weakly in
$X\widehat\otimes\ell^2$ and strongly in
$L^2_{\mathrm{loc}}(\mathbb R^N;\ell^2)$.
\end{lemma}
\begin{proof}
Fix an orthonormal basis $(\phi_k)$ of the scalar Hilbert space $X$.
Write its internal coefficient vectors as
$c_{nk}=\langle\phi_k,U_n\rangle_X\in\ell^2$.
Gram--Schmidt the first $n$ vectors and extend the resulting finite
orthonormal list to an internal orthonormal basis. A unitary $W_n$
can therefore be chosen so that
\[
 W_nc_{nk}\in\operatorname{span}\{e_1,\ldots,e_k\},\qquad k\le n.
\]
For each fixed $k$, these vectors have a convergent subsequence in a
fixed finite-dimensional space. Diagonal extraction gives strong
convergence of every coefficient vector. Weak compactness gives a
weak limit of $W_nU_n$, with exactly those coefficients.

For $\chi\in C_c^\infty$, the scalar multiplication map
$K_\chi:X\to L^2$, $f\mapsto\chi f$, is compact. If $P_L$ is the
scalar projection onto the first $L$ basis vectors, then
\[
 \|K_\chi(I-P_L)\|_{X\to L^2}\longrightarrow0.
\]
Tensoring with the identity on $\ell^2$ preserves this norm.
The transformed $P_LU_n$ converge strongly because there are only
$L$ coefficient vectors; after multiplication by $\chi$, the tails
are uniformly small. This proves local strong convergence.
\end{proof}

\begin{proof}[Compactness assertion of Theorem~\ref{thm:r11-vectorcompact}]
The localized Sobolev inequality has constants independent of internal
dimension: use the Hilbert-valued Bessel-potential representation and
its norm triangle inequality, then spatial localization and the sum of
component quadratic norms. Thus local mass vanishing for an
$H^s$-bounded sequence implies vanishing in $L^p(\mathcal H)$ for
$2<p<2N/(N-2s)$, or every finite $p>2$ if $2s\ge N$.
In the local case take $p=2\kappa$; in the Hartree case use
$p=4N/(2N-m)$ and Hardy--Littlewood--Sobolev. Hence
$\mathcal F(U_n)\to\kappa>0$ excludes local mass vanishing.

Translate a nonzero amount of local mass into a fixed ball and use
Lemma~\ref{lem:r11-internalcompact}. We obtain a nonzero weak limit
$U$, strong local $L^2$ convergence, and by interpolation strong local
convergence in the interaction exponent. Put $R_n=U_n-U$.
The quadratic quantities split. For local powers the Brezis--Lieb
argument for $\|U\|_{\mathcal H}^{2\kappa}$ gives
\[
 \mathcal F(U_n)=\mathcal F(U)+\mathcal F(R_n)+o(1).
\]
For Hartree, the term $\operatorname{Re}\langle U,R_n\rangle$
tends to zero in $L^{p/2}$ by local convergence and the small
$L^p$ tail of $U$. Also
$\mathcal D(\rho_U,\rho_{R_n})\to0$: approximate $U$ in
$L^p(\mathcal H)$ by a compactly supported smooth finite-dimensional
field. Its potential is bounded and decays at infinity, while $R_n$
tends locally to zero in $L^2$. HLS controls the approximation error.
This proves the same splitting for Hartree.

Let $M_1>0,M_2$ and $A_1,A_2$ be the limiting contributions.
Then $M_1+M_2=M_c$, $A_1+A_2=1$, and sharp interpolation yields
\[
 M_c^{\kappa-1}\le
 M_1^{\kappa-1}A_1+M_2^{\kappa-1}A_2.
\]
Both masses cannot be positive, since each would be strictly below
$M_c$. Thus $M_2=0$, and the inequality then forces $A_2=0$.
The convergence is strong in $H^s$. The limit is an optimizer,
already proved to be completely polarized.
\end{proof}

\subsection{Phase coercivity and the internal symmetry kernel}
For completeness, the scalar phase fact used below follows directly
from the density norm. Normalize a real optimizer by $M(Q)=M_c$,
$A(Q)=1$, and put $\omega_*=(\kappa-1)/M_c$. The form
\[
 L_Q=a(D)+\omega_*-V_{Q^2}
\]
is nonnegative: differentiate the sharp inequality at
$Q+i\sqrt t\,h$. If $h$ belongs to its kernel, write
$x=A(h)$ and $y=M(h)/M_c$. The supporting-functional inequality
for the density norm, sharp interpolation for $h$, and weighted
arithmetic--geometric mean give
\[
 x+(\kappa-1)y
 \le\kappa x^{1/\kappa}y^{(\kappa-1)/\kappa}
 \le x+(\kappa-1)y.
\]
Equality implies $x=y$ and $h^2=yQ^2$ by strict convexity.
The kernel of a nonnegative form is linear; apply the same conclusion
to $h+Q$ to get $h=cQ$. The potential is compact as a form
$H^s\to H^{-s}$, by local Rellich and H\"older in the local case,
or the compact-potential approximation and HLS in the Hartree case.
Consequently
\begin{equation}\label{eq:r11-vectorcoercivity}
 \langle h,L_Qh\rangle\ge c_Q\|h\|_{H^s}^2,
 \qquad \langle h,Q\rangle_{L^2}=0.
\end{equation}
This holds for complex $h$ by real--imaginary splitting, and it can
be summed over arbitrarily many components with the same constant.

At $Qe_1$, write a perturbation as $(h+ik)e_1+Z$ with $Z\perp e_1$.
The Hessian of $E+\omega_*M/2$ is
\[
 \langle h,L_+h\rangle+\langle k,L_Qk\rangle
 +\sum_{j\ge2}\bigl(
  \langle\operatorname{Re}Z_j,L_Q\operatorname{Re}Z_j\rangle
 +\langle\operatorname{Im}Z_j,L_Q\operatorname{Im}Z_j\rangle\bigr).
\]
Its kernel in the phase and internal transverse sector is exactly
\[
 \{Q\eta:\ \eta\in\mathcal H,
                 \operatorname{Re}\langle\eta,e_1\rangle=0\}.
\]
For $\mathcal H=\mathbb C^d$, this is the $(2d-1)$-dimensional real
tangent space to the internal unitary orbit. Any additional full
Hessian kernel lies in the scalar real sector $L_+$; no assumption
on that sector is made.

\subsection{Quantitative polarization without real-sector nondegeneracy}
The preceding compactness and phase coercivity also give a quantitative
statement for arbitrary fields, not just stationary or traveling ones.
Define the scale-invariant interpolation deficit
\[
 \mathfrak d(U)=1-
 \frac{\mathcal F(U)}{\kappa(M(U)/M_c)^{\kappa-1}A(U)}
\]
and the squared distance to a completely polarized field by
\[
 \mathfrak p(U)=\inf_{\substack{\|\zeta\|_{\mathcal H}=1\\q\in H^s}}
 \left\{\frac{M(U-\zeta q)}{M(U)}
             +\frac{A(U-\zeta q)}{A(U)}\right\}.
\]
Both are defined for every nonzero $U\in H^s(\mathcal H)$.

\begin{theorem}[Dimension-independent quantitative polarization]
\label{thm:r11-quantitativepolarization}
There is $C<\infty$, depending only on the fixed scalar dispersion
and interaction, such that
\begin{equation}\label{eq:r11-polarizationdeficit}
 \mathfrak p(U)\le C\mathfrak d(U)
\end{equation}
for every nonzero Hilbert-valued field $U$ and every internal dimension.
No nondegeneracy or uniqueness of the scalar real-sector ground state
is required.
\end{theorem}
\begin{proof}
Amplitude and spatial scaling preserve both quantities, so impose
$M(U)=M_c$, $A(U)=1$. Write
$\|U\|_X^2=M(U)/M_c+A(U)$.
Since $\mathfrak p(U)\le2$, failure of a uniform constant would
produce $U_n$ with $\mathfrak d(U_n)\to0$ and
$\mathfrak d(U_n)/\mathfrak p(U_n)\to0$.
After the symmetries of Theorem~\ref{thm:r11-vectorcompact},
$U_n\to Qe_1$ strongly in $H^s$.
Align $\int QU_n$ with the positive real first internal direction
and write $U_n=(u_n,Z_n)$. Each component of $Z_n$ is complex
$L^2$-orthogonal to $Q$, and $u_n\to Q$, $Z_n\to0$ in $H^s$.

The interaction expansion in these transverse directions is
\begin{equation}\label{eq:r11-transverseexpansion}
 \mathcal F(u_n,Z_n)=\mathcal F(u_n)
 +\kappa\langle V_{|u_n|^2},\|Z_n\|^2\rangle
 +o(\|Z_n\|_X^2).
\end{equation}
For local powers, set $a=|u_n|^2$ and $b=\|Z_n\|^2$ pointwise.
When $1<\kappa\le2$, the nonnegative Taylor remainder of
$(a+b)^\kappa$ after its linear term in $b$ is at most $C_\kappa b^\kappa$.
When $\kappa>2$, it is at most
$C_\kappa(a^{\kappa-2}b^2+b^\kappa)$.
H\"older and the dimension-independent $L^{2\kappa}$ Sobolev bound
make these respectively $O(\|Z_n\|_X^{2\kappa})$ and
$O(\|Z_n\|_X^4+\|Z_n\|_X^{2\kappa})$, since $u_n$ stays bounded.
Both are $o(\|Z_n\|_X^2)$.
For Hartree the remainder is exactly
$\mathcal D(\|Z_n\|^2,\|Z_n\|^2)=O(\|Z_n\|_X^4)$.

Put $m_n=M(Z_n)$ and $a_n=A(Z_n)$. Scalar sharp interpolation gives
\[
 \mathcal F(u_n)\le
 \kappa(1-m_n/M_c)^{\kappa-1}(1-a_n).
\]
Combining its expansion with \eqref{eq:r11-transverseexpansion}
and potential form convergence yields
\begin{align*}
 \kappa-\mathcal F(U_n)
 &\ge\kappa\left[
 a_n+\frac{\kappa-1}{M_c}m_n
       -\langle V_{Q^2},\|Z_n\|^2\rangle\right]
       -o(\|Z_n\|_X^2)\\
 &=\kappa\sum_j\langle Z_{nj},L_QZ_{nj}\rangle
       -o(\|Z_n\|_X^2)\\
 &\ge c\|Z_n\|_X^2.
\end{align*}
The constant is independent of component number. Taking the candidate
$e_1u_n$ in the distance gives
$\mathfrak p(U_n)\le\|Z_n\|_X^2$.
As $\mathfrak d(U_n)=(\kappa-\mathcal F(U_n))/\kappa$,
this contradicts the assumed ratio and proves the theorem.
\end{proof}

The theorem controls internal polarization rather than the full
scalar shape. It therefore remains valid even when scalar optimizer
uniqueness or real-sector quantitative stability is unavailable.

